\subsection{Integer central restriction and cubic constraints}
\begin{proposition}[Integer central necessary system]
Let $\mathcal C$ be pseudo-unitary with commutative Grothendieck ring,
$D=\operatorname{FPdim}(\mathcal C)$, characters $\chi_E$, and formal
codegrees $f_E$. Put $\beta_E=D/f_E$, and let $A_E$ denote the distinguished
simple of $Z(\mathcal C)$ belonging to $E$. A complete finite restriction
column domain gives a finite system of integer equations with bounded
nonnegative integer variables that includes restriction Gram identities,
distinguished $S$-symmetry, duality, and
\[
 N_{A_E,A_F}^{A_G}
 =\sum_Z\frac{\beta_E\beta_F\beta_{G^*}
 \chi_E(FZ)\chi_F(FZ)\chi_{G^*}(FZ)}{D^2\operatorname{FPdim}(Z)}.
\]
Every categorification supplies an integer solution of this system.
\end{proposition}
\begin{proof}
By \cite[Lemma~4.3 and Corollary~4.5]{DeshpandeCenter},
$S_{A_E,Z}=\beta_E\chi_E(FZ)$ in the unnormalized modular convention.
The displayed identity is the Verlinde formula for the center of dimension
$D^2$. Introduce a nonnegative integer selection variable for each allowed
restriction of each distinguished simple, with sum one for each character.
The tensor unit is fixed separately; its character is located by equality
with the positive dimension character. High variables count actual central
simples per complete restriction-duality orbit. Gram, modular symmetry and
low-duality equations are linear in these variables. Expanding the displayed
formula in the full character field makes every nonconstant coordinate zero
and its constant coordinate a nonnegative integer auxiliary variable.
Diagonal Gram bounds and the sum of positive constant contributions give
finite upper bounds. Clearing denominators produces the stated system.
No integrality is imposed on relative twist-trace weights.
\end{proof}
The eigenvalues $\beta_E\chi_E(FZ)/\operatorname{FPdim}(Z)$ of the integer
fusion matrix of $A_E$ also have equal multiplicities throughout each full
Galois orbit. These yield further exact linear equations. Integer interval
and gcd propagation are necessary tests. A covector annihilating all columns
modulo an integer $m>1$, but not the right-hand side, certifies nonmembership
in the column lattice; rational annihilators are checked in the same manner.
Signed integer lattice membership alone is weaker than a nonnegative
bounded solution. Numerical optimization only proposes vectors, which are
checked against every exact original equation and bound. Any negative
pseudo-unitary result still requires a complete signed and dimension-character
argument before it is promoted to a general-complex exclusion.
The actual center of $\mathrm{Fib}\boxtimes\mathrm{Vec}_{C_4}$ provides a
64-simple positive integer control with complex character values.
The dimension identity further gives the exact bound
\[
 N_{A_E,A_F}^{A_G}\leq
 \left\lfloor\frac{\beta_E\beta_F}{\beta_G}\right\rfloor.
\]
If a candidate eigenvalue has a Galois conjugate absent from the complete
remaining domain, its nonnegative multiplicity is zero; this rule may be
iterated, retaining the fixed constituents. Row content division, exact
proportional-row identification, and zero-row deletion preserve the integer
solution set. Their signed row scales also lift any compact-system
annihilator back to a covector checked on every original column.
